\documentclass[10pt,a4paper]{article}
\usepackage[T1]{fontenc}
\usepackage[utf8]{inputenc}
\usepackage[french]{babel}
\usepackage{newtxtext}
\usepackage{geometry}
\usepackage{xcolor}
\usepackage{enumitem}
\usepackage{tabularx}
\usepackage{array}
\usepackage{fancyhdr}
\usepackage{hyperref}
\usepackage{microtype}
\geometry{left=1.8cm,right=1.8cm,top=1.85cm,bottom=1.8cm}
\definecolor{vert}{HTML}{176C47}
\definecolor{gris}{HTML}{51575C}
\hypersetup{colorlinks=true,urlcolor=vert}
\setlength{\parindent}{0pt}
\setlength{\parskip}{3pt}
\setlist[itemize]{leftmargin=*,topsep=2pt,itemsep=0pt}
\pagestyle{fancy}
\fancyhf{}
\lhead{\textcolor{vert}{SportsVision AI}}
\rhead{Mohamed Amine Bensserrar}
\cfoot{\thepage}
\renewcommand{\headrulewidth}{0.3pt}
\newcommand{\partie}[1]{\vspace{6pt}{\large\bfseries\color{vert}#1}\par\vspace{2pt}\hrule\vspace{4pt}}
\newcommand{\etape}[3]{\textbf{#1. #2}\par #3\par\vspace{3pt}}
\newcommand{\qa}[2]{\textbf{Q. #1}\par \textbf{R.} #2\par\vspace{4pt}}
\begin{document}
{\LARGE\bfseries Comprendre mon projet SportsVision AI}\par
{\color{gris}Fiche de révision simple -- français niveau B1 -- pas un texte à réciter}\par
\textbf{Étudiant :} Mohamed Amine Bensserrar \qquad
\textbf{Sujet :} classification automatique de photos sportives
\vspace{4pt}

\partie{1. Quel est le but du projet ?}
Nous avons créé une application web qui reconnaît le sport présent dans une photo. Elle classe les images en trois catégories : \textbf{basketball, football et tennis}. L'utilisateur peut envoyer une photo, voir le sport prédit et le pourcentage de confiance. Il peut aussi envoyer plusieurs photos et télécharger un fichier ZIP avec trois dossiers.

\textbf{Pourquoi ?} Le but est d'automatiser le classement des photos. Le programme reconnaît des éléments visuels : ballon, terrain, raquette, joueurs, etc.

\partie{2. Le travail réalisé, étape par étape}
\etape{1}{Préparer Google Colab}{Nous avons utilisé Python, TensorFlow et Keras dans Google Colab pour entraîner un modèle de Deep Learning. Colab permet d'exécuter le notebook sans installer tout le travail sur notre PC.}
\etape{2}{Télécharger le dataset}{Avec \texttt{kagglehub}, nous avons téléchargé le dataset public \emph{Sports Classification} de Kaggle. Le dataset avait déjà les dossiers \texttt{train}, \texttt{valid} et \texttt{test}. Nous avons gardé seulement trois sports pour notre projet.}
\etape{3}{Vérifier les données}{Nous avons compté les images de chaque sport. Les dossiers donnent les noms des classes (les étiquettes ou \emph{labels}).}
\begin{center}
\begin{tabular}{lrrr}
\hline
\textbf{Sport} & \textbf{Train} & \textbf{Validation} & \textbf{Test}\\\hline
Basketball &169&5&5\\
Football &191&5&5\\
Tennis &131&5&5\\\hline
\textbf{Total}&\textbf{491}&\textbf{15}&\textbf{15}\\\hline
\end{tabular}
\end{center}
\textbf{Train :} apprendre. \quad \textbf{Validation :} suivre l'apprentissage. \quad \textbf{Test :} vérifier le résultat final.

\etape{4}{Charger les images}{Nous avons utilisé \texttt{image\_dataset\_from\_directory} pour lire les photos. Toutes les images passent à \textbf{224 x 224 pixels}, en RGB (trois couleurs). Le \textbf{batch size = 16} signifie que le modèle traite 16 images par groupe. Les images de train sont mélangées (\texttt{shuffle=True}) ; celles de validation et de test ne sont pas mélangées.}
\etape{5}{Choisir MobileNetV2}{C'est un réseau CNN déjà entraîné sur ImageNet. Il sait reconnaître des caractéristiques générales des images. Nous avons utilisé le \textbf{Transfer Learning} pour réutiliser ces connaissances. Cela évite de créer un gros modèle à partir de zéro.}
\etape{6}{Bloquer les anciennes couches}{Avec \texttt{base\_model.trainable=False}, nous avons bloqué les poids de MobileNetV2. Ces poids ne changent pas. Nous entraînons seulement notre nouvelle couche de classification.}
\etape{7}{Construire notre modèle}{Pipeline réel : \textbf{image 224 x 224} $\to$ \textbf{Rescaling} $\to$ \textbf{MobileNetV2} $\to$ \textbf{GlobalAveragePooling2D} $\to$ \textbf{Dense(3) + Softmax}.}
\begin{itemize}
\item \textbf{Rescaling :} change les pixels de 0--255 vers -1--1, pour MobileNetV2.
\item \textbf{MobileNetV2 :} extrait les caractéristiques utiles de la photo.
\item \textbf{GlobalAveragePooling2D :} résume les cartes de caractéristiques en 1280 valeurs.
\item \textbf{Dense(3) :} donne un résultat pour chacun des trois sports.
\item \textbf{Softmax :} transforme les résultats en trois probabilités dont la somme est 100 \%.
\end{itemize}
\etape{8}{Configurer l'apprentissage}{Nous avons choisi \textbf{Adam} (optimiseur), \textbf{learning rate = 0,001}, \textbf{sparse categorical crossentropy} (fonction de perte) et \textbf{accuracy} (mesure des bonnes réponses).}
\etape{9}{Entraîner le modèle}{Nous avons utilisé \texttt{model.fit(..., epochs=10)}. Une \textbf{époque} est un passage sur toutes les 491 images de train. Pendant ces dix époques, seule la couche finale a appris. La validation est utilisée pour observer les résultats.}
\etape{10}{Évaluer avec le test}{Nous avons utilisé \texttt{model.evaluate(test\_ds)}. Sur les \textbf{15 images de test}, le modèle a trouvé \textbf{15 bonnes réponses} : accuracy 100 \%, test loss environ 0,0215. \textbf{Attention :} 15 images, c'est trop peu pour garantir 100 \% sur toutes les nouvelles photos.}
\etape{11}{Construire la matrice de confusion}{Nous avons comparé les vraies étiquettes et les classes prédites. La matrice est : basketball \textbf{5/5}, football \textbf{5/5}, tennis \textbf{5/5}. Sur ce petit test, precision, recall et F1-score valent aussi 100 \% pour chaque classe.}
\etape{12}{Sauvegarder les fichiers}{Nous avons exporté \texttt{sports\_model.keras} (réseau entraîné) et \texttt{labels.json} (noms et ordre des sports : basketball, football, tennis). L'ordre est important pour lire les probabilités.}
\etape{13}{Créer le backend}{Avec \textbf{FastAPI} en Python, nous avons créé une API. Elle charge le modèle sauvegardé au démarrage. Pour une image, \texttt{POST /api/predict} vérifie le fichier, le transforme en RGB 224 x 224, exécute le modèle et renvoie le sport et les probabilités. Le backend ne réentraîne pas le modèle.}
\etape{14}{Créer le frontend}{Nous avons utilisé \textbf{HTML, CSS et JavaScript} pour permettre à l'utilisateur d'importer une image, voir son aperçu et lire la prédiction. JavaScript envoie l'image à FastAPI, puis affiche la réponse.}
\etape{15}{Ajouter le classement automatique}{La fonctionnalité multiple accepte jusqu'à \textbf{20 photos}. \texttt{POST /api/organize} renvoie les catégories, les compteurs et les résultats. \texttt{POST /api/organize/zip} crée le ZIP avec les dossiers \texttt{basketball/}, \texttt{football/} et \texttt{tennis/}. Les photos originales sont gardées dans l'archive ; le serveur ne les conserve pas de façon permanente.}
\etape{16}{Tester et documenter}{Nous avons exécuté des tests backend : \textbf{5 réussis, 0 échec}. Les tests navigateur ont aussi réussi, sans erreur JavaScript. Les captures, le notebook, le code GitHub et le rapport documentent le travail.}

\partie{3. Comment une prédiction fonctionne ?}
\textbf{Exemple :} je choisis une photo de football.
\begin{enumerate}[leftmargin=*,itemsep=0pt]
\item JavaScript envoie la photo à \texttt{/api/predict}.
\item FastAPI contrôle le format et lit les pixels.
\item La photo est convertie en RGB et redimensionnée en 224 x 224.
\item Le modèle applique lui-même le \texttt{Rescaling}. Le serveur ne le fait pas une deuxième fois.
\item MobileNetV2 extrait les informations ; la couche Dense + Softmax calcule les trois probabilités.
\item Le backend prend la plus grande probabilité et renvoie \texttt{predicted\_sport}, \texttt{confidence} et \texttt{probabilities}.
\item Le site montre le sport et les pourcentages.
\end{enumerate}

\partie{4. Concepts essentiels en mots simples}
\begin{tabularx}{\textwidth}{>{\bfseries}p{3.6cm}X}
CNN & Réseau adapté aux images : il apprend les formes et les textures.\\
Convolution / kernel & Petit filtre qui passe sur l'image pour chercher des informations.\\
Feature map & Résultat d'un filtre : une carte de caractéristiques.\\
ReLU & Fonction qui garde les valeurs positives et met les valeurs négatives à zéro.\\
Pooling & Réduit les dimensions des informations.\\
Transfer Learning & Réutiliser un modèle qui a déjà appris sur beaucoup d'images.\\
Poids gelés & Anciennes valeurs du modèle qui ne changent pas pendant notre entraînement.\\
Loss & Mesure de l'erreur. Plus elle est basse, mieux c'est en général.\\
Accuracy & Pourcentage de réponses correctes.\\
Époque & Un tour complet sur les données d'entraînement.\\
Batch & Petit groupe d'images traité ensemble.\\
Softmax & Donne une probabilité pour chaque classe.\\
Matrice de confusion & Tableau des bonnes et mauvaises prédictions pour chaque classe.\\
Overfitting & Modèle qui apprend trop les exemples et généralise mal.\\
\end{tabularx}

\partie{5. Questions possibles du professeur : réponses directes}
\qa{Quel est le problème traité ?}{Reconnaître automatiquement le sport dans une photo et organiser les images par catégorie.}
\qa{Pourquoi trois sports seulement ?}{Pour avoir un projet clair et entraîner un premier modèle fonctionnel avec trois classes.}
\qa{D'où viennent les images ?}{Du dataset public \emph{Sports Classification} sur Kaggle.}
\qa{Pourquoi trois dossiers train, validation et test ?}{Train sert à apprendre, validation à suivre les résultats, test à vérifier le modèle à la fin.}
\qa{Combien d'images avez-vous utilisées ?}{491 pour train, 15 pour validation et 15 pour test, soit 521 images au total.}
\qa{Pourquoi 224 x 224 ?}{Parce que c'est la taille d'entrée choisie pour notre MobileNetV2.}
\qa{Pourquoi RGB ?}{Parce que notre modèle attend des images avec trois canaux de couleur.}
\qa{C'est quoi un CNN ?}{Un réseau de neurones qui apprend des informations visuelles, comme les formes et les textures.}
\qa{Quel est le rôle d'un kernel ?}{C'est un petit filtre qui parcourt l'image pour détecter des caractéristiques.}
\qa{C'est quoi une feature map ?}{C'est la carte produite par un filtre après une convolution.}
\qa{Pourquoi MobileNetV2 ?}{Il est relativement léger et déjà entraîné. Il permet d'apprendre plus vite avec peu d'images.}
\qa{C'est quoi le Transfer Learning ?}{On garde les connaissances d'un modèle pré-entraîné et on apprend une nouvelle tâche.}
\qa{Pourquoi avez-vous gelé MobileNetV2 ?}{Pour garder ses connaissances et entraîner seulement notre nouvelle couche finale.}
\qa{Que fait GlobalAveragePooling2D ?}{Elle réduit les cartes de caractéristiques à un vecteur de 1280 valeurs.}
\qa{Pourquoi Dense(3) ?}{Parce que nous avons trois sports à reconnaître.}
\qa{Pourquoi Softmax ?}{Pour obtenir trois probabilités et choisir la plus grande.}
\qa{Pourquoi la normalisation ?}{MobileNetV2 attend des valeurs adaptées. Notre couche Rescaling transforme les pixels de 0--255 vers -1--1.}
\qa{Pourquoi Adam ?}{C'est un optimiseur qui aide le modèle à ajuster les poids pendant l'apprentissage.}
\qa{Pourquoi cette fonction de perte ?}{Sparse categorical crossentropy convient à nos trois classes avec des étiquettes entières.}
\qa{C'est quoi une époque ?}{Un passage sur toutes les images d'entraînement. Nous avons fait dix époques.}
\qa{C'est quoi batch size 16 ?}{Le modèle traite les images par groupes de 16.}
\qa{Pourquoi shuffle=True pour train ?}{Pour varier l'ordre des images pendant l'apprentissage.}
\qa{Quelle est la différence entre accuracy et loss ?}{Accuracy mesure les bonnes réponses ; loss mesure l'erreur du modèle.}
\qa{Que montre votre matrice de confusion ?}{Cinq bonnes prédictions pour chacun des trois sports et aucune erreur sur les 15 images de test.}
\qa{Votre modèle a-t-il vraiment 100 \% de précision ?}{Sur nos 15 images de test, oui. Mais ce test est petit : cela ne prouve pas 100 \% pour toutes les images.}
\qa{Pourquoi sauvegarder labels.json ?}{Pour associer correctement chaque sortie du modèle au nom du sport.}
\qa{Quel est le rôle de FastAPI ?}{Recevoir les images, utiliser le modèle et renvoyer les résultats au site.}
\qa{Quel est le rôle de JavaScript ?}{Envoyer les fichiers à l'API et afficher la réponse dans la page.}
\qa{Pourquoi avoir ajouté le ZIP ?}{Pour classer automatiquement plusieurs photos en trois dossiers et les télécharger facilement.}
\qa{Avez-vous réentraîné le modèle dans le site web ?}{Non. Le site charge le fichier .keras déjà entraîné dans Colab.}
\qa{Quelles sont les limites du projet ?}{Le test est petit, il y a seulement trois sports et le modèle peut se tromper sur des images nouvelles.}
\qa{Quelle amélioration proposeriez-vous ?}{Ajouter plus de données, plus de sports et un plus grand jeu de test.}
\qa{Comment savez-vous que le site fonctionne ?}{Nous avons testé les prédictions, le classement par lots, les erreurs et le ZIP. Les tests ont réussi.}

\partie{6. Les sept idées à retenir}
\begin{enumerate}[leftmargin=*,itemsep=1pt]
\item \textbf{But :} reconnaître et ranger des photos sportives.
\item \textbf{Données :} Kaggle, 3 sports, train/validation/test = 491/15/15.
\item \textbf{Modèle :} MobileNetV2 + Transfer Learning ; ancienne base gelée.
\item \textbf{Apprentissage :} 224 x 224, batch 16, Adam, 10 époques.
\item \textbf{Résultat :} 15/15 sur le petit test, sans garantie générale.
\item \textbf{Site :} FastAPI + HTML/CSS/JS + modèle .keras.
\item \textbf{Bonus utile :} plusieurs photos $\to$ catégories $\to$ ZIP.
\end{enumerate}
\medskip
{\small\color{gris}Sources du projet : \url{https://github.com/aminebenserrar/sportsvision-ai} et notebook Colab \url{https://colab.research.google.com/drive/1zgeVvl2XZNKWtXm44oMiwRtb2vpyfGwb}.}
\end{document}
